\section{Composable excursions with genuine old-wake
readback}\label{composable-excursions-with-genuine-old-wake-readback}

\subsection{Convention, standard word, and finite mature
restrictions}\label{returns-2}

Use the unchanged HC contract: complete finite black set B, integer head
p, headings N,E,S,W=0,1,2,3. Read white/right or black/left, toggle,
then move. State t is before update t. Rotations are clockwise
quarter-turns, never an unannounced reflection of the LR rule.

Use exactly the standard word w and thirteen-cell NE seed K from
\passthrough{\lstinline!verification/guarded\_scattering/check.py!}. Let
P\_i,H\_i be its 104 pre-update positions and headings, with P\_0=0,
H\_0=E and d=(2,2). Its minimum future value of x+y is zero. Write w\_a
for its cyclic word beginning at phase a, and R\_k for clockwise
rotation through k quarter-turns. The return checker pins the earlier
source hash; the paper reproduction replays both the 30-rule library and
the return module.

Put

\begin{lstlisting}
s(x,y)=x+y,  r(x,y)=x-y,  ell(x,y)=y-x,  G=-160.
\end{lstlisting}

For a phase j, orientation k, and a linear form f with f(R\_k
d)\textgreater0, define

\begin{lstlisting}
M(j,k,f,G) = XOR_(i=0..103) XOR_(L_i<=a<=u_i)
                  { R_k(P_i-P_j)+a R_k d },
L_i=ceil((G-f(R_k(P_i-P_j)))/f(R_k d)),
u_i=0 if i<j, and -1 otherwise.                         (2.1)
\end{lstlisting}

This is finite. It is the exact parity of all planned standard writes
before a phase-j collision, restricted to f\textgreater=G. It is a
description of a finite restriction, not an infinite physical input. A
finite approach realizes it once the old physical endpoint and the
synthetic past omitted before that approach lie outside the comparison
half-plane. Those obligations are proved below, not assumed from
observed repetition.

Our first collision interface has head 0, heading H\_57=S and field

\begin{lstlisting}
M1 union {0,q2(n),q3} union H,
M1=M(57,0,s,G),             H subset {s<G},             (2.2)
\end{lstlisting}

where H is any finite complete historical black field. The two-marker
version omits q3. Future packets add further markers above the current
packet's entire read extent; their noninterference is proved in
Subsection \ref{returns-8}.

\subsection{Finite A and C blocks, and the first long
flight}\label{returns-3}

\subsubsection{3.1 First turn A}\label{first-turn-a}

The exact 832-update orbit from (M1 union \{0\},0,H\_57) ends at

\begin{lstlisting}
v_A=(-12,12), h_A=S,
\end{lstlisting}

with 580 black cells and a certified standard NW tail w\_58. Every A
read has -10\textless=s\textless=9 and every tail read has
s\textgreater=G. Both infinite checkers accept this complete reference
endpoint. Its initial field has 484 cells.

Let Q\_B,i be the phase positions of w\_58 from this endpoint. Then

\begin{lstlisting}
q2(n)=Q_B,14+n*(-2,2),
\end{lstlisting}

and the first reference read there is exactly 104n+14 updates later. The
finite ray calculation enumerates every i and integer z with
q2(0)-Q\_B,i=z*(-2,2). All n+z are nonnegative for n\textgreater=100,
and min(104z+i)=14. The expected letter is R. This checks first contact,
not just that the proposed marker eventually lies on the reference
route.

The marker q2(n) is strictly left of all A reads, for every
n\textgreater=100. The fixed q3 is absent from A and from every future
NW phase ray. The old field H is outside s\textgreater=G; the A prefix
and its NW reference remain inside. First-difference coupling proves the
actual A block and the subsequent 104n+14-update flight B, ending at the
black marker q2. Its absolute packet time is 104n+846.

\subsubsection{3.2 Distant reversal C}\label{distant-reversal-c}

Let M2=M(72,3,ell,G). The exact 51272-update orbit from

\begin{lstlisting}
(M2 union {0},0,(H_72+3) mod 4)
\end{lstlisting}

ends at v\_C=(-24,-44), heading N, with 2569 black cells and a certified
SE reference tail w\_18. Its initial field has 485 black cells. All its
finite reads have ell\textgreater=G and, after translation by q2(n),
s\textgreater=G. Its bounds are -75\textless=s\textless=63 and
-76\textless=r\textless=54 in LOCAL coordinates.

The true field before this C block agrees with this kernel on

\begin{lstlisting}
Omega_C(n)={s>=G and ell(v-q2(n))>=G}.                  (3.1)
\end{lstlisting}

Indeed every black cell of the A endpoint and q3 lies strictly below the
ell boundary. So does every point of the synthetic standard past
preceding the A endpoint. Their maxima are finite; subtracting
ell(q2(n))=4n+26 makes each required strict inequality an integer affine
inequality in n.~They hold at n=100 with nonnegative slope thereafter.
Hence inside (3.1), the actual B writes equal exactly the mature M2
writes from (2.1). The marker adds \{0\}. The arbitrary H contributes
nothing because of the s boundary.

Thus C is an actual literal block on the complete input. No all-future
conclusion for the TRUE orbit is inferred from its isolated-kernel tail:
the older field outside Omega\_C will matter during the reverse flight.

\subsection{Exact first return to old memory, not an assumed invisible
archive}\label{returns-4}

At C's end the head is

\begin{lstlisting}
q2(n)+v_C=(-2n-39,2n-33), h=N,                        (4.1)
\end{lstlisting}

at packet time 104n+52118. Its isolated reference follows w\_18
southeast. Let Q\_D,i(0) be the 104 phase positions from (-39,-33),
heading N. For general n the phase positions are Q\_D,i(0)-n*(2,-2).

The true field and this reference differ by the following exact XOR:

\begin{lstlisting}
A_endpoint XOR B_writes XOR (q2(n)+M2) XOR {q3} XOR H. (4.2)
\end{lstlisting}

The finite C writes cancel because both trajectories execute them
identically. The two copies of the marker q2 also cancel. This identity
retains every older write; it does not identify the true state with a
cleaned reference.

All reverse reference phase rays have s values strictly below both the
NW B write stream and q2(n)+M2, while still having s\textgreater=G.
Therefore they cannot read those two terms or H. The remaining
first-difference query is against the fixed finite A endpoint and q3.

For each phase i and black point b in A\_endpoint union \{q3\},
enumerate integers z satisfying

\begin{lstlisting}
b-Q_D,i(0)=z*(2,-2).
\end{lstlisting}

There are 158 phase/cell candidates. Every n+z is nonnegative for
n\textgreater=100. The minimum of 104z+i is -170, attained at phase
index 38 and b=q*=(-38,-32). No candidate belongs to q3. The reference
letter there is R and the actual cell is black. Hence the actual reverse
flight has EXACT length

\begin{lstlisting}
104n-170 = 104(n-2)+38,                                (4.3)
\end{lstlisting}

and ends in the old black read at packet time 208n+51948. Before it, all
turns agree with the SE reference. At it, the real turn is L rather than
R.

\subsubsection{The black cell was created by the same actual
trajectory}\label{the-black-cell-was-created-by-the-same-actual-trajectory}

In the incoming NE plan relative to phase 57, solve

\begin{lstlisting}
q*=P_i-P_57+a*d,
104a+i-57<0.
\end{lstlisting}

The last such read is i=89,a=-17, at time -1736, and w\_89=R. The A
block never reads q*; B is disjoint by its s range; C misses the return
region below; and (4.3) proves that D does not read q* earlier. Thus its
black lifetime is exactly

\begin{lstlisting}
(208n+51948)-(-1736)=208n+53684.                       (4.4)
\end{lstlisting}

In every train realization below, the incoming flight has more than 1736
updates, so this birth occurs AFTER the finite input time, even in the
first packet. It is not an initially supplied marker, an intrinsic
pre-input lifetime, or a cell that is simply presumed eventually
relevant.

\subsection{A constant full comparison field after an arbitrarily long
excursion}\label{returns-5}

Use the fixed region

\begin{lstlisting}
Omega={s>=G and r>=G}.                                (5.1)
\end{lstlisting}

At the old read q*, the true field restricted to Omega is independent of
n. This is the closure-enabling fact. The arbitrary H is absent by its s
guard. The distant q2 marker and ALL C reads are outside Omega for every
n\textgreater=100: their r values are -4n-26 plus a bounded local value,
strictly below G. This does not say they were erased; they remain in the
complete true endpoint.

For completeness, here is an explicit finite XOR formula for the field
in Omega. Start with A\_endpoint union \{q3\}. For each Q\_B,i put

\begin{lstlisting}
h_i=floor((r(Q_B,i)-G)/4).
\end{lstlisting}

Add by XOR all Q\_B,i+j*(-2,2) with 0\textless=j\textless=h\_i and
s(Q\_B,i)\textgreater=G. All h\_i=100. The terminal 14 B reads at epoch
n are outside Omega. Thus these are exactly the B writes inside Omega,
not a truncation by simulation time.

For the reverse flight write its complete epochs as v=j-n, so
-n\textless=v\textless=-3. For each Q\_D,i(0) put

\begin{lstlisting}
l_i=ceil((G-r(Q_D,i(0)))/4).
\end{lstlisting}

Add by XOR Q\_D,i(0)+v*(2,-2) for l\_i\textless=v\textless=-3 when
s(Q\_D,i(0))\textgreater=G. All l\_i\textgreater=-n for
n\textgreater=100. Finally add the first 38 phase positions at v=-2, if
they lie in Omega. Restrict the resulting finite field to Omega.

The resulting field C\_old has 1446 cells. Its exact hash and the 3618 B
terms and 3893 D terms are recorded in the report. This is equality of
the ENTIRE inside field, including the old A scar and both relevant
write streams. It is not equality on the next few read positions or a
bounded mask inferred from matching poses. Every lower/upper epoch
omission is justified by the linear forms above.

\subsubsection{5.1 Two-marker output}\label{two-marker-output}

From (C\_old minus \{q3\},q*,S), an exact 9308-update block ends at
(-22,-60),S, with a certified SE w\_14 tail. All these finite reads and
the infinite tail stay in Omega. First-difference induction with (5.1)
proves that the complete two-marker state is permanently standard by
packet time 208n+61256.

This is already an all-n positive family with a real unbounded-age
readback. Its whole-input realizations from K have FIFTEEN initial black
cells when there is no extra debris. No uniform onset bound depending
only on this cardinality is implied; coordinates still encode n.

\subsubsection{5.2 Reset to the original NE
interface}\label{reset-to-the-original-ne-interface}

The marker q3 is not read in that 9308-update recovery. It is the
phase-15 first-visit cell exactly 10401 updates along the SE reference,
and is read white/R there by the unmarked reference. The third initial
black cell changes that turn. The complete finite block from
(C\_old,q*,S) has 23661 updates and ends at

\begin{lstlisting}
v_out=(203,-244), E, standard NE phase 94.             (5.2)
\end{lstlisting}

All its reads and its infinite NE tail stay in Omega. Two infinite
criteria accept its complete 3453-black comparison endpoint. The same
coupling proves the actual THREE-marker packet permanently standard by
208n+75609 when no later packet interrupts it.

The initial q3 has not been manufactured by this process: omitting it
gives the different, permanent SE future just proved. This explicit
distinction prevents treating the reset as an autonomous instruction
generator.

\subsection{Single-packet theorem and protection under later
interruptions}\label{returns-6}

\textbf{Theorem 1.} Every complete state (2.2), for n\textgreater=100
and arbitrary finite H subset \{s\textless G\}, has the three-marker
behavior just described and reaches NE standard phase 94 by time
208n+75609. The two-marker state has the SE output from Subsection
\ref{returns-5}.1. All supplied black history remains physically
present.

\textbf{Proof.} Subsections \ref{returns-3} and \ref{returns-4} give the exact A, B, C, D execution and
first old read. Subsection \ref{returns-5} gives the exact whole-field
restriction and protected finite E continuation and tail. The comparison
paths remain within the regions on which field agreement was
established, so a first divergent read is impossible. This is a
finite-stage induction plus an infinite-tail certificate, not
extrapolation from the sampled n values. \(\square\)

For the three-marker packet, every read BEFORE (5.2) satisfies

\begin{lstlisting}
-109 <= s <= 59.                                      (6.1)
\end{lstlisting}

The upper bound also covers all initial M1 cells and the three markers.
The NW and SE flights have constant s on each phase ray; their
arbitrarily long lengths do not enlarge (6.1). The finite A, C, E blocks
supply the other extrema. At the C boundary, (4.1) gives x=-2n-39, so no
n-independent backward x-displacement or bounded-diameter packet
footprint can replace this argument.

A later marker may interrupt the final NE tail. The single-packet
theorem alone would not guarantee old-history protection after that
interruption. The next section supplies the necessary nested-guard
induction for the FULL composition.

\subsection{Finite arithmetic connector and marker
ordering}\label{returns-7}

A packet outputs pose v\_out, orientation NE, phase 94. After
m\textgreater=100 full periods and phase adjustment to 57, its next
primary marker is displaced by

\begin{lstlisting}
Delta(m)=v_out+(m+1)*d+P_57-P_94
        =(2m+208,2m-242).                             (7.1)
\end{lstlisting}

Thus \(s(\Delta(m))=4m-34\ge366\). The next diagonal guard starts
at \(s(\Delta(m))+G\ge206\), while all previous packet black cells
and reads not in inherited history have s\textless=59. The minimum guard
slack is 147, strictly positive. Every synthetic pre-output standard
write is also below the next boundary; this is checked for each phase at
its latest possible pre-output epoch, with all earlier epochs decreasing
s further.

It follows that, at the next primary, the full inside field consists
exactly of its M1, its three still-unread markers, and still later
preinstalled markers. The actual NE flight writes and the mature
synthetic writes agree there because the omitted pre-flight past is
entirely outside. Inherited earlier history is below an even older guard
and hence below the new one.

The companion markers of this next packet cannot interrupt that NE
approach: q2(n) and q3 have r respectively -4n-26 and 437 relative to
the primary, outside ALL two-sided incoming NE phase lines. This is a
geometric disjointness test, not merely checking a finite flight. Nor
can markers from a still later packet cause a premature contact. Their
minimum relative s is -83; (7.1) places them strictly above both the
current packet's entire read maximum and the maximum s of an incoming
flight before its primary. The latter is at most primary s+1, as checked
from the 104 phases with their latest pre-collision epochs.

These inequalities prove the chronological order of all planned
contacts, not just compatibility of consecutive endpoint poses. For
m\textgreater100 their left sides increase by 4(m-100); no further
trajectory simulation is required.

\subsection{Arbitrarily many genuine returns from one complete finite
input}\label{returns-8}

Choose ANY integer k\textgreater=0 and ANY integer lists

\begin{lstlisting}
m_1,...,m_k>=100,    n_1,...,n_k>=100.
\end{lstlisting}

Start with pose p\_0=0, phase a\_0=0, orientation NE. For i=1,\ldots,k
set

\begin{lstlisting}
Q_i=p_(i-1)+(m_i+1_[57<a_(i-1)])*d+P_57-P_(a_(i-1));
p_i=Q_i+v_out,     a_i=94.                            (8.1)
\end{lstlisting}

Install at INPUT time all three cells

\begin{lstlisting}
Q_i,  Q_i+q2(n_i),  Q_i+q3.
\end{lstlisting}

The complete input is K plus these cells plus any finite black G0 subset
\{s\textless0\}; its ant starts at 0,E. All these sets are disjoint, and
the input has exactly 13+3k+\textbar G0\textbar{} black cells. Markers
are never installed during execution. For k=0 the initial state is
already on the certified NE future.

\textbf{Theorem 2 (closed arbitrary-return trains).} Every such finite
input follows the prescribed packets with all past writes retained and
reaches permanent standard P104 by

\begin{lstlisting}
T = sum_(i=1..k) [104m_i+(57-a_(i-1) mod 104)
                            +208n_i+75609].            (8.2)
\end{lstlisting}

The output after k\textgreater0 packets is p\_k, heading E, NE phase 94.
Each packet contains the genuine post-input-created black lifetime
(4.4).

\textbf{Proof.} The canonical initial ray has s\textgreater=0, so it
cannot read G0. The first primary guard is beyond every K cell and the
entire synthetic pre-start past; at m\_1=100 its slack is at least 243.
Its companion markers miss the incoming NE rail, and later markers are
beyond the finite read envelope. Therefore the first collision has the
correct complete restricted field (2.2).

Apply Theorem 1 up to its output, but keep all later markers in the
physical state. Subsection \ref{returns-7} proves that they have not
altered this prefix, that the next collision has exactly the required
restricted field, and that its guard is strictly above all preceding
history. Induct on i. The guards strictly increase in the SAME absolute
linear form s. Between collisions the previous certified tail remains
above its previous guard; within each new packet all reads stay above
the new guard. Consequently no later reset reactivates a retired region.
After the last packet there is no remaining forward perturbation and the
certified NE tail persists. Summing the exact block lengths proves
(8.2). \(\square\)

The assertion about arbitrary finite G0 is a proved initial guard, not a
deletion of an unknown future-active set. Complete-state train replays
retain G0 and every black point from every earlier excursion.

\subsubsection{Small input and repeated-return
examples}\label{small-input-and-repeated-return-examples}

With m\_1=n\_1=100 and no G0, the three markers are

\begin{lstlisting}
(204,205), (-11,416), (381,-55).
\end{lstlisting}

Together with K this is a SIXTEEN-black finite seed. It writes (166,173)
at time 8721 and next reads it black at time 83205, an exact lifetime of
74484 updates. It has a certified NE output at time 106866, at
(407,-39), with 7796 black cells.

The retained six-packet example has m\_i=100 and
\(n_i=(100,101,150,103,211,100)\). Its THIRTY-ONE-black input executes six
such returns and has a certified NE output at time 675566, at
(2447,-249), with 50671 black cells. These are examples of the all-k
proof, not a claim about all sixteen- or thirty-one-black initial
states.

\subsection{A finite-stage rank for return packets}\label{returns-9}

The packet phases are an exact segmentation of a restricted true orbit,
NOT PG secondary gates or intrinsic S\_D reads. Use a finite-stage
lexicographic rank

\begin{lstlisting}
(packets not yet completed, remaining stage index, local fuel).
\end{lstlisting}

The stage order is incoming NE flight, A, B, C, D, E. Each flight has
the exact remaining update count from the formulas above. In each finite
A, C or E block, choose its finite read rectangle, INCLUDING the last
endpoint, and use the scaled integer r=4 clock already proved in
clock.md. The unchanged exterior contribution cancels. Its computable
upper bound over all internal fields and poses yields the nonnegative
integer fuel

\begin{lstlisting}
ceil((U_Q-C_Q(X))/epsilon_Q),
\end{lstlisting}

which decreases at every internal literal update. This remains valid for
the TRUE complete field: unread differences inside that rectangle merely
contribute constant terms, and the universal clock inequality holds for
every colouring. Transitions to the next stage decrease the middle
coordinate, allowing fuel to reset. Completing a packet decreases the
first coordinate and starts the next incoming flight; completing the
last packet halts the certificate program. The actual ant continues
forever along its highway.

Unbounded excursion n makes the two flight counters large but does not
create new unbudgeted stages: the exact XOR return query and full
return-region equality proved that there is exactly the certified old
collision and recovery. That is the closure fact on which the rank
relies. The clock did not prove the guard, the return-field equality, or
all-input coverage.

This rank carries the finite packet description. It is not a new
all-state physical rank, a global PG/S\_D event bound, or a proof that
arbitrary finite seeds enter this interface.
